The integration of Artificial Intelligence (AI) into healthcare has reshaped Clinical Decision Support Systems (CDSS), precision medicine, and prognostic modeling, where adoption depends not only on predictive performance but on explainability, transparency, and trust, especially in safety-critical settings where opaque model behavior undermines accountability and patient care \cite{ghassemi2021falsehope,holzinger2019causability}. Meanwhile, privacy regulations such as HIPAA \cite{hipaa1996} and GDPR \cite{gdpr2016} restrict access to Electronic Health Records (EHR), fragmenting data across institutional silos and limiting both model training and validation. This tension has positioned Synthetic Data Generation (SDG) as a promising enabler of \rev{privacy-conscious} analysis and digital patient-twin construction \cite{goncalves2020generation,azizi2021can}; but to support explainable AI, synthetic data must itself be generated by mechanisms that preserve clinically meaningful structure.

Yet existing synthetic tabular generators fall short. Biomedical data are heterogeneous, continuous physiological measurements alongside discrete diagnostic and procedural codes, with non-linear dependencies and severe class imbalance. Generative Adversarial Networks (GANs) such as CTGAN \cite{xu2019modeling} and MedGAN \cite{choi2017generating} suffer mode collapse, underrepresenting rare but clinically critical subpopulations; diffusion models \cite{kotelnikov2023tabddpm} inject stochastic noise that can obscure clinically meaningful thresholds and correlations. Both operate as black boxes, giving little insight into which clinical constraints are satisfied or violated and making errors hard to trace to specific features, so they hinder rather than support trustworthy pipelines.

Large Language Models (LLMs) have emerged as flexible generators of serialized tabular data (GReaT \cite{borisov2022language}, RealTabFormer \cite{solatorio2023realtabformer}), capturing global distributions and rare categorical values better than prior architectures through pre-trained representations. Their process, however, remains opaque: unconstrained LLMs can hallucinate medically implausible records, a pregnancy diagnosis for a male patient, an invalid ICD code, with no signal indicating why. Post-hoc explanations on downstream models cannot compensate for synthetic data that encodes hidden inconsistencies, so explainability must be addressed during generation, not retrofitted afterward.

\noindent \textbf{Problem Statement.}
Formally, biomedical tabular synthesis requires approximating the joint distribution of heterogeneous Electronic Health Record (EHR) features while satisfying multiple non-commensurate constraints at distinct granularities: (i) \emph{syntactic validity} of values (e.g., legal ICD formats and numeric ranges), (ii) \emph{local semantic consistency} among related attributes (e.g., demographic--diagnosis compatibility), (iii) \emph{feature-level dependency preservation} across correlated variables, and (iv) \emph{global row-level coherence} of realistic patient profiles. Medical data compound this with extreme class imbalance, where rare but critical conditions occupy a vanishing fraction of the distribution. Existing paradigms are ill-suited: likelihood objectives optimize marginal token probabilities without structured medical logic; adversarial and diffusion models collapse rare modes or corrupt correlations; and reinforcement-learning approaches rely on scalar or textual feedback that cannot localize errors across semantic levels. The unresolved problem is thus how to optimize a high-dimensional generative policy under \emph{multi-objective, hierarchical clinical constraints} whose violations are sparse, heterogeneous, and not expressible as a single differentiable loss.



% To bridge the gap between statistical plausibility and clinical validity, we introduce \textbf{DISCO-TAB} (DIScriminator-guided COntrol for TABular synthesis). Unlike prior approaches that rely on singular scalar feedback or unsupervised likelihood maximization, DISCO-TAB orchestrates a fine-tuned LLM with a Hierarchical Reinforcement Learning framework. We employ a multi-objective discriminator system that evaluates synthesis at four distinct granularities: \textit{token-level} (syntax validity), \textit{sentence-level} (semantic coherence), \textit{feature-level} (statistical correlation), and \textit{row-level} (global clinical logic). Furthermore, to address the pervasive issue of class imbalance in medical datasets, we introduce Inverse-Frequency Reward Shaping (IFRS), a mechanism that dynamically scales RL rewards to incentivize the accurate generation of rare minority classes.



\noindent\textbf{Proposed Approach.}
To bridge the gap between statistical plausibility and clinical validity, we introduce \textbf{DISCO-TAB} (DIScriminator-guided COntrol for TABular synthesis), a unified framework that explicitly operationalizes biomedical tabular synthesis as a constrained, sequential decision-making problem. Rather than relying on a single likelihood objective or monolithic discriminator, DISCO-TAB couples a fine-tuned Large Language Model (LLM) with a hierarchical reinforcement learning (RL) optimization strategy, enabling the generator to receive \rev{structured, multi-granular} feedback aligned with clinical reasoning. Specifically, we design a multi-objective discriminator ensemble that evaluates generated records at four complementary granularities: \emph{token-level} constraints enforce syntactic correctness of numerical values and categorical codes; \emph{sentence-level} constraints penalize local semantic contradictions (e.g., demographic--diagnosis inconsistencies); \emph{feature-level} constraints preserve statistically significant clinical dependencies automatically discovered from real EHR data; and \emph{row-level} constraints enforce global coherence of entire patient profiles. 

To further address the endemic class imbalance of biomedical datasets, we introduce \emph{Inverse-Frequency Reward Shaping} (IFRS), which dynamically scales reinforcement signals to incentivize faithful generation of rare but clinically critical patient subpopulations. This hierarchical, reward-shaped RL formulation allows DISCO-TAB to localize and correct violations at the appropriate semantic level during generation, rather than filtering errors post hoc. As a result, the framework improves minority-class coverage, preserves fine-grained feature correlations, and suppresses medically implausible hallucinations. Across a 23-dataset benchmark spanning clinical, financial, chemical, and web domains, DISCO-TAB attains top-tier downstream utility among generators with full benchmark coverage and, unlike several diffusion and language-model baselines, consistently avoids single-class collapse on the small, imbalanced clinical datasets, while leading the full-coverage field on correlation preservation and keeping disclosure-risk indicators calibrated.
These results demonstrate that explicitly encoding clinical constraints into the optimization loop is essential for trustworthy biomedical data synthesis, particularly in small-sample and high-imbalance regimes where existing generative models fail.


\noindent\textbf{Our contributions are as follows.} The individual building blocks we employ (autoregressive LLM generators, Proximal Policy Optimization (PPO)-based reinforcement learning, discriminators, and correlation-based association measures) are not independently novel. The central contribution of DISCO-TAB is their integration into a hierarchical, multi-granular feedback architecture for clinical tabular synthesis, together with two mechanisms tailored to biomedical data (automated dependency discovery and inverse-frequency reward shaping), and its empirical evaluation. Concretely:

1. \textit{Hierarchical Clinical Verification:} We propose a multi-granular discriminator system that enforces validity from the atomic token level (e.g., valid lab values) to the holistic record level, coupling these heterogeneous validity signals with reinforcement-learning-based generator optimization to reduce medical hallucinations, a form of supervision that a single adversarial or likelihood objective cannot express.

2. \textit{Automated Constraint Discovery:} We introduce an automated pre-processing module that extracts latent statistical dependencies from raw EHR data, initializing feature-specific discriminators without requiring manual rule engineering by clinicians. \rev{This discovery is currently based on pairwise associations (see Section~\ref{sec:discussion} for the resulting limitation).}

3. \textit{Large-Scale Empirical Validation:} We evaluate DISCO-TAB against thirteen baselines across 23 datasets and 21 metrics spanning six evaluation dimensions, with repeated runs. DISCO-TAB attains top-tier downstream utility and the strongest correlation-preservation fidelity among generators with full benchmark coverage, and consistently avoids single-class collapse. We also state plainly the datasets and metrics on which competing methods lead (Section~\ref{sec:experiments}).
